\chapter{Linear Time-Invariant Systems}

So far we have spent our energy on \emph{signals} -- objects that carry
information through time. A microphone trace, an accelerometer stream, the EMG
crackle from a forearm muscle: these are all signals. But signals rarely sit
still. We feed them through things. A microphone signal passes through a tone
control. An accelerometer reading passes through a smoothing filter before we
trust it. The raw EMG from your wrist passes through a rectifier and an envelope
detector before any gesture is recognized.

Each of those ``things'' is a \emph{system}: a box that takes a signal in and
produces a signal out. This chapter is about a remarkably tractable and
ubiquitous class of systems -- the \textbf{linear
time-invariant} (LTI) systems. The whole edifice of classical signal processing
-- filters, the frequency response, convolution, the transfer function -- rests
on two innocent-looking assumptions. We will see that those two assumptions are
enough to pin a system down \emph{completely}: tell me how an LTI
system reacts to a single tap, and I can tell you how it reacts to \emph{any}
input whatsoever.

\begin{focusbox}
A \textbf{system} is a map $T$ that turns an input signal $x$ into an output
signal $y = T\{x\}$. We will write this as
\[
  x[n] \;\longrightarrow\; \boxed{\;T\;}\; \longrightarrow\; y[n]
\]
for discrete time, or $x(t)\to y(t)$ for continuous time. The chapter's central
claim: if $T$ is linear and time-invariant, then it is fully described by one
signal -- its response to an impulse -- and it acts on every sinusoid simply by
scaling and phase-shifting it.
\end{focusbox}

We will work mostly in discrete time, because that is what runs on the
microcontroller strapped to the wrist. Everything has a continuous-time twin;
we point those out as we go.

\section{Systems: the input--output view}
\label{ch12:system-view}

Forget, for a moment, what is inside the box. A system is defined entirely by
its behavior: which output it assigns to each input. This is the \emph{black
box} or \emph{input--output} view, and it is the right altitude from which to
start.

\begin{example}[Some systems, with their rules]
\label{ch12:zoo}
Each of these maps an input sequence $x[n]$ to an output sequence $y[n]$.
\begin{itemize}
  \item \textbf{Amplifier:} $y[n] = 3\,x[n]$. Scale by a constant.
  \item \textbf{Delay:} $y[n] = x[n-2]$. Shift two samples to the right.
  \item \textbf{Two-point averager:} $y[n] = \tfrac12\big(x[n] + x[n-1]\big)$.
  \item \textbf{Accumulator:} $y[n] = \sum_{k=-\infty}^{n} x[k]$ (a running sum).
  \item \textbf{Squarer:} $y[n] = x[n]^2$.
  \item \textbf{Rectifier:} $y[n] = \abs{x[n]}$ (the heart of EMG envelope
        detection).
  \item \textbf{Time-warp:} $y[n] = x[2n]$ (decimation / downsampling).
\end{itemize}
\end{example}

These look superficially similar -- each is ``some formula for $y$ in terms of
$x$.'' But they split sharply into two camps once we ask the two questions that
organize this entire chapter:

\begin{enumerate}
  \item Is the system \textbf{linear}? (Does it respect scaling and adding?)
  \item Is it \textbf{time-invariant}? (Does it treat all moments alike?)
\end{enumerate}

The amplifier, delay, and averager pass both tests. The squarer and rectifier
fail linearity. The time-warp fails time-invariance. By the end of the chapter
you will see exactly why those failures matter -- and why the systems that pass
both tests are the ones we can analyze with Fourier methods.

\begin{intuitionbox}
Think of a system as a \emph{machine that reshapes signals}. A filter that
removes hiss reshapes the signal by deleting its high-frequency content. A
reverb reshapes by smearing energy forward in time. The two LTI properties are
exactly the conditions under which this reshaping can be understood, frequency by
frequency, without the frequencies interfering with one another.
\end{intuitionbox}

\section{Linearity: scaling and superposition}
\label{ch12:linearity}

The first property is the one you already know from linear algebra, transplanted
to signals. Recall (from the companion linear-algebra book) that a map $T$ is
\emph{linear} when it commutes with the two vector-space operations: scaling and
addition.

\begin{definition}[Linearity]
\label{ch12:def-linear}
A system $T$ is \textbf{linear} if for all inputs $x_1, x_2$ and all scalars
$a, b$ (real or complex),
\[
  T\{\,a\,x_1[n] + b\,x_2[n]\,\} \;=\; a\,T\{x_1[n]\} + b\,T\{x_2[n]\}.
\]
\end{definition}

This single equation bundles two ideas:
\begin{itemize}
  \item \textbf{Homogeneity (scaling):} doubling the input doubles the output.
        $T\{a\,x\} = a\,T\{x\}$.
  \item \textbf{Additivity (superposition):} the response to a sum is the sum of
        the responses. $T\{x_1 + x_2\} = T\{x_1\} + T\{x_2\}$.
\end{itemize}

\begin{connectionbox}
In the linear-algebra book a linear map satisfied $T(a\vec u + b\vec v) = a
T\vec u + b T\vec v$ and was represented by a \emph{matrix}. A discrete signal
\emph{is} a vector (Chapter~3 treated $x[n]$ as a vector indexed by $n$), and a
linear system \emph{is} a (possibly infinite) matrix acting on it:
$y = \mat{H}\,x$. Time-invariance, coming next, is the extra condition that
makes this matrix have a beautiful repeating structure -- and that structure is
what convolution is.
\end{connectionbox}

\begin{whybox}
Why obsess over superposition? Because it is the engine of \emph{decompose and
reassemble}. If I can write a complicated input as a sum of simple pieces,
\[
  x = \sum_k c_k\,\phi_k,
\]
and I know what the system does to each simple piece $\phi_k$, then linearity
hands me the answer for free:
\[
  y = T\{x\} = \sum_k c_k\,T\{\phi_k\}.
\]
The entire Fourier program is exactly this: choose the $\phi_k$ to be sinusoids,
because (we will prove) the system treats each sinusoid in the most trivial
possible way. Linearity is the permission slip that lets us add the answers back
up.
\end{whybox}

\subsection*{Testing for linearity}

To \emph{prove} a system linear you must verify the definition in general. To
\emph{disprove} it, a single counterexample suffices.

\begin{example}[The amplifier is linear]
$T\{x\}[n] = 3x[n]$. Then
$T\{ax_1 + bx_2\}[n] = 3(ax_1[n] + bx_2[n]) = a\,(3x_1[n]) + b\,(3x_2[n])
= a\,T\{x_1\}[n] + b\,T\{x_2\}[n]$. Linear. \checkmark
\end{example}

\begin{example}[The squarer is not linear]
\label{ch12:squarer-nonlinear}
$T\{x\}[n] = x[n]^2$. Try scaling: $T\{2x\}[n] = (2x[n])^2 = 4x[n]^2$, but
$2\,T\{x\}[n] = 2x[n]^2$. Since $4x[n]^2 \neq 2x[n]^2$ in general, homogeneity
fails. The squarer is nonlinear.
\end{example}

\begin{trapbox}
``Linear'' does \emph{not} mean ``the graph is a straight line,'' and it does
\emph{not} mean ``contains no constant term.'' Consider $y[n] = x[n] + 5$. Its
input--output graph is a straight line, yet it is \emph{not} linear:
$T\{0\} = 5 \neq 0$, but every linear system must send the zero signal to the
zero signal (set $a=b=0$ in Definition~\ref{ch12:def-linear}). Such ``line plus
offset'' maps are called \emph{affine}, not linear. The offset breaks
superposition: $T\{x_1+x_2\} = x_1+x_2+5$, while $T\{x_1\}+T\{x_2\} =
x_1+x_2+10$.
\end{trapbox}

\section{Time invariance: the system has no clock}
\label{ch12:ti}

The second property says the system does not care \emph{when} you run the
experiment. Shift the input in time, and the output shifts by the same amount,
unchanged in shape.

\begin{definition}[Time invariance]
\label{ch12:def-ti}
A system $T$ is \textbf{time-invariant} if, whenever $y[n] = T\{x[n]\}$, then for
every integer shift $n_0$,
\[
  T\{\,x[n - n_0]\,\} \;=\; y[n - n_0].
\]
In words: delaying the input by $n_0$ delays the output by $n_0$ and does nothing
else.
\end{definition}

\begin{intuitionbox}
A time-invariant system has \emph{no internal clock and no calendar}. It does not
know whether it is Monday or whether the signal arrived early or late. Run the
same input today or tomorrow and you get the same output, just slid along the
time axis to match. A microphone preamp is time-invariant. A system whose gain
ramps up over the first second after power-on is \emph{not} -- it knows how long
it has been running.
\end{intuitionbox}

\subsection*{Testing for time invariance}

The recipe is mechanical and worth doing carefully, because the most common
student error lives right here.

\begin{walkbox}
\textbf{Recipe.} To test time invariance, compute two things and compare:
\begin{enumerate}
  \item \textbf{Shift-then-process.} Let $x_{\text{sh}}[n] = x[n-n_0]$, push it
        through the system, get $T\{x_{\text{sh}}\}[n]$.
  \item \textbf{Process-then-shift.} Take the original output
        $y[n] = T\{x\}[n]$ and shift it: $y[n-n_0]$.
\end{enumerate}
If the two agree \emph{for all} $x$ and all $n_0$, the system is time-invariant.
The trap is in step 1: wherever an $n$ appears in the system's formula
\emph{attached to $x$}, you must replace $x[\,\cdot\,]$ by the shifted input,
but any \emph{free} $n$ (an explicit dependence on the current time) stays put --
and that mismatch is exactly what breaks time invariance.
\end{walkbox}

\begin{example}[The delay is time-invariant]
$T\{x\}[n] = x[n-2]$, so $y[n] = x[n-2]$.
\emph{Shift-then-process:} feeding $x[n-n_0]$ gives output $x[(n-2)-n_0] =
x[n-2-n_0]$.
\emph{Process-then-shift:} $y[n-n_0] = x[(n-n_0)-2] = x[n-2-n_0]$.
They match. Time-invariant. \checkmark
\end{example}

\begin{example}[The downsampler is not time-invariant]
\label{ch12:downsampler}
$T\{x\}[n] = x[2n]$, so $y[n] = x[2n]$.
\emph{Shift-then-process:} feed $x[n-n_0]$; the output is $x[2n - n_0]$.
\emph{Process-then-shift:} $y[n-n_0] = x[2(n-n_0)] = x[2n - 2n_0]$.
These differ ($x[2n-n_0]$ vs.\ $x[2n-2n_0]$) for, e.g., $n_0=1$. So
downsampling is \emph{not} time-invariant: a one-sample input delay does not
produce a one-sample output delay. (This is why resampling, Chapter~22, needs
its own careful theory.)
\end{example}

\begin{example}[A multiplier by $n$ is not time-invariant]
$T\{x\}[n] = n\,x[n]$. Here the free $n$ out front is the culprit.
Shift-then-process: $n\,x[n-n_0]$. Process-then-shift: $(n-n_0)\,x[n-n_0]$.
The factors $n$ and $(n-n_0)$ disagree -- the system weights later samples more
heavily, so it manifestly has a clock.
\end{example}

\section{Putting them together: LTI}
\label{ch12:lti}

\begin{definition}[LTI system]
\label{ch12:def-lti}
A system is \textbf{LTI} if it is both linear (Definition~\ref{ch12:def-linear})
and time-invariant (Definition~\ref{ch12:def-ti}).
\end{definition}

Of the zoo in Example~\ref{ch12:zoo}: the amplifier, delay, two-point averager,
and accumulator are LTI; the squarer and rectifier fail linearity; the
downsampler fails time invariance. The averager and accumulator deserve a second
look because they are the prototypes of the two great families of filters we
build later (FIR and IIR).

\begin{example}[The two-point averager is LTI]
$y[n] = \tfrac12(x[n] + x[n-1])$. Linearity follows because it is a sum of two
linear pieces (scaled, delayed copies of $x$), and any sum of linear operations
is linear. Time invariance follows because each piece -- a scaling and a fixed
delay -- is time-invariant, and the sum of time-invariant operations is
time-invariant. This ``sum of scaled delays'' template is the giveaway: any such
combination is automatically LTI, and \emph{every} LTI system turns out to be
a (possibly infinite) sum of scaled delays. We prove that general result over the
next two sections.
\end{example}

\section{The impulse: a single tap}
\label{ch12:impulse}

To unlock LTI systems we need one special signal: the \textbf{unit impulse}
(also called the unit sample, or Kronecker delta in discrete time),
\[
  \delta[n] = \begin{cases} 1, & n = 0,\\ 0, & n \neq 0. \end{cases}
\]
It is a single spike of height $1$ at time zero and silence everywhere else --
the cleanest possible ``tap'' you can give a system.

\begin{figure}[H]
\centering
\begin{tikzpicture}
\begin{axis}[
  width=10cm, height=4.2cm,
  xlabel={$n$}, ylabel={$\delta[n]$},
  xmin=-4.5, xmax=4.5, ymin=-0.2, ymax=1.25,
  xtick={-4,-3,-2,-1,0,1,2,3,4}, ytick={0,1},
  axis lines=left, ymajorgrids, grid style={gray!20},
  every axis plot/.append style={sigblue, thick},
]
\addplot+[ycomb, mark=*, mark options={fill=sigblue}]
  coordinates {(-4,0)(-3,0)(-2,0)(-1,0)(0,1)(1,0)(2,0)(3,0)(4,0)};
\end{axis}
\end{tikzpicture}
\caption{The unit impulse $\delta[n]$: a single sample of height one at $n=0$.}
\label{ch12:fig-impulse}
\end{figure}

The reason the impulse is special is a \emph{sifting} identity that says any
signal is a sum of shifted, scaled impulses.

\begin{formalbox}
\textbf{The sifting (sampling) representation.} Every discrete signal can be
written as
\[
  x[n] \;=\; \sum_{k=-\infty}^{\infty} x[k]\,\delta[n-k].
\]
\emph{Why this identity holds:} fix an index $n$. On the right-hand side, the term
$\delta[n-k]$ is zero unless $k=n$, in which case it equals $1$. So the entire
sum collapses to the single survivor $x[n]\cdot 1 = x[n]$. The sum reconstructs
the signal sample by sample.
\end{formalbox}

That equation looks like bookkeeping, but read it as a \emph{decomposition}: the
signal $x$ is a weighted superposition of building blocks $\delta[n-k]$, where
the weight on the block sitting at time $k$ is precisely the sample value $x[k]$.
We have decomposed $x$ into impulses. We are now exactly one application of
linearity away from the whole theory.

\begin{historybox}
The formal concept of a Linear Time-Invariant system arose during the expansion of transcontinental
telephony and feedback amplifier design at Bell Laboratories in the 1920s and 1930s.
Engineers like Harold Black (who conceived the negative-feedback amplifier while riding the ferry
across the Hudson River in 1927), Harry Nyquist, and Hendrik Bode needed mathematical principles to
guarantee that cascading thousands of vacuum-tube repeaters along phone cables would not spontaneously
erupt into howling oscillation. Restricting attention to systems satisfying superposition and time-invariance
simplified the analysis: the entire system behavior could be predicted from its response to elementary
impulses.

The continuous-time impulse $\delta(t)$ was introduced by Paul Dirac in the late 1920s as an idealized
``unit mass concentrated at a point.'' Classical mathematicians protested that Dirac's delta violated
function theory, until Laurent Schwartz (1950) founded the theory of distributions. In discrete time,
the Kronecker delta $\delta[n]$ is completely free of distribution-theoretic paradoxes: it is simply
a sequence that equals $1$ at $n=0$ and $0$ elsewhere. Because any discrete signal can be decomposed into
a weighted train of shifted impulses, an LTI system's response to any input is completely determined by
its impulse response $h[n]$.

Linearity and time-invariance are of course physical idealizations. Real amplifiers clip when overdriven,
semiconductors drift with ambient temperature, and mechanical transducers wear down. But within their
linear operating ranges, LTI systems grant us total predictive power. Chapter~13 (\emph{Convolution})
turns this impulse-decomposition insight into a precise computational engine, unpacking the geometry
and mechanics of the convolution sum.
\end{historybox}

\section{The impulse response and the convolution sum}
\label{ch12:impulse-response}

Now feed the impulse into an LTI system and watch what comes out. Give it a
name.

\begin{definition}[Impulse response]
\label{ch12:def-h}
The \textbf{impulse response} of a system $T$ is its output when the input is the
unit impulse:
\[
  h[n] \;=\; T\{\delta[n]\}.
\]
\end{definition}

The impulse response is a signal -- the system's ``signature,'' the shape of its
reaction to a single tap. A bell's impulse response is its ringing decay. A
room's impulse response is its reverberation (clap once, record what comes back).
Here is the theorem that makes $h[n]$ the most important object in the chapter.

\begin{theorem}[An LTI system is its impulse response]
\label{ch12:thm-conv}
If $T$ is LTI with impulse response $h[n] = T\{\delta[n]\}$, then for
\emph{every} input $x[n]$ the output is the \textbf{convolution}
\[
  y[n] \;=\; (x \conv h)[n] \;=\; \sum_{k=-\infty}^{\infty} x[k]\,h[n-k].
\]
\end{theorem}

\begin{whybox}
\textbf{The proof is just our two properties, in order.} Start from the sifting
representation and apply the system:
\[
  y[n] = T\{x[n]\} = T\!\left\{ \sum_{k} x[k]\,\delta[n-k] \right\}.
\]
\emph{Step 1 -- linearity.} The system commutes with the (weighted) sum. The
weights $x[k]$ are constants as far as $n$ is concerned, so they pull out:
\[
  y[n] = \sum_{k} x[k]\; T\{\delta[n-k]\}.
\]
\emph{Step 2 -- time invariance.} The system's response to the \emph{shifted}
impulse $\delta[n-k]$ is just the shifted impulse response,
$T\{\delta[n-k]\} = h[n-k]$. Substitute:
\[
  y[n] = \sum_{k} x[k]\,h[n-k] = (x\conv h)[n]. \qquad\blacksquare
\]
This two-step deduction connects the entire behavior of the system to its impulse response: linearity breaks the input into a sum of weighted impulses, while time invariance ensures the response to each shifted impulse is a shifted copy of one master signature $h$.
\end{whybox}

\begin{happeningbox}
What is really happening: an LTI system \emph{stamps a copy of its impulse
response at every input sample}, scaled by that sample's value, and adds all the
stamped copies up. Picture a drummer hitting a ringing cymbal. Each hit at time
$k$ with strength $x[k]$ launches one cymbal-ring $x[k]\,h[n-k]$. The sound you
hear at time $n$ is the sum of all the still-ringing decays from every earlier
hit. Convolution is ``overlapping echoes, added.''
\end{happeningbox}

This gives immense predictive leverage. A system can be infinitely complicated
as a rule, but if it is LTI then \emph{one experiment} -- measure $h[n]$ --
predicts its response to every conceivable input. We measure $h$, we store it,
and convolution does the rest.

\begin{figure}[H]
\centering
\begin{tikzpicture}
\begin{axis}[
  width=12.5cm, height=4.5cm,
  title={Convolution as stamped, scaled, added echoes},
  title style={font=\small},
  xlabel={$n$}, ymin=-0.2, ymax=1.15,
  xmin=-0.6, xmax=10.6, xtick={0,2,4,6,8,10},
  ytick=\empty, axis lines=left, legend pos=north east,
  legend style={font=\footnotesize, draw=gray!40},
]
\addplot+[ycomb, mark=*, sigblue, mark options={fill=sigblue}, forget plot]
  coordinates {(0,0.9)(1,0.72)(2,0.576)(3,0.461)(4,0.369)};
\addlegendimage{sigblue, thick}\addlegendentry{copy at $k{=}0$, scale $0.9$}
\addplot+[ycomb, mark=triangle*, targetgreen, mark options={fill=targetgreen}, forget plot]
  coordinates {(3,0.5)(4,0.4)(5,0.32)(6,0.256)(7,0.205)};
\addlegendimage{targetgreen, thick}\addlegendentry{copy at $k{=}3$, scale $0.5$}
\addplot+[ycomb, mark=square*, quotientorange, mark options={fill=quotientorange}, forget plot]
  coordinates {(6,0.3)(7,0.24)(8,0.192)(9,0.154)(10,0.123)};
\addlegendimage{quotientorange, thick}\addlegendentry{copy at $k{=}6$, scale $0.3$}
\end{axis}
\end{tikzpicture}
\caption{Each input sample launches a scaled, shifted copy of the impulse
response $h[n] = 0.8^{\,n}u[n]$. The output is their pointwise sum. Three input
samples are shown ($x[0]{=}0.9$, $x[3]{=}0.5$, $x[6]{=}0.3$); a real signal
launches one such echo per sample.}
\label{ch12:fig-conv-echoes}
\end{figure}

\begin{marginnotebox}
Convolution gets a whole chapter of its own (Chapter~13): the flip-and-slide
picture, its commutativity ($x\conv h = h\conv x$), and fast algorithms. Here we
only need the \emph{fact} that LTI output $=$ input convolved with $h$. We will
lean on it immediately.
\end{marginnotebox}

\begin{example}[Convolving by hand]
\label{ch12:hand-conv}
Let the impulse response be the two-point averager, $h[n] = \tfrac12\delta[n] +
\tfrac12\delta[n-1]$, i.e.\ $h = (\tfrac12, \tfrac12)$ at $n=0,1$. Let the input
be the short pulse $x = (1, 3, 2)$ at $n=0,1,2$. Then
\[
  y[n] = \sum_k x[k]\,h[n-k]
       = \tfrac12 x[n] + \tfrac12 x[n-1].
\]
March $n$ forward:
\[
\begin{aligned}
 y[0] &= \tfrac12(1) + \tfrac12(0) = 0.5,\\
 y[1] &= \tfrac12(3) + \tfrac12(1) = 2.0,\\
 y[2] &= \tfrac12(2) + \tfrac12(3) = 2.5,\\
 y[3] &= \tfrac12(0) + \tfrac12(2) = 1.0.
\end{aligned}
\]
So $y = (0.5, 2.0, 2.5, 1.0)$. Note the output is one sample \emph{longer} than
the input: convolving a length-$L$ signal with a length-$M$ impulse response
gives length $L+M-1$. The averaging visibly smooths the jumpy input.
\end{example}

\section{FIR and IIR: finite versus infinite memory}
\label{ch12:fir-iir}

The \emph{length} of $h[n]$ tells you how long the system remembers a tap.

\begin{itemize}
  \item If $h[n]$ has only finitely many nonzero samples, the system has
        \textbf{Finite Impulse Response (FIR)}. The averager
        $h=(\tfrac12,\tfrac12)$ is FIR: a tap is forgotten after one sample.
  \item If $h[n]$ rings forever, the system has \textbf{Infinite Impulse
        Response (IIR)}. The accumulator's impulse response is the unit step
        $h[n]=u[n]=1$ for $n\ge0$ -- it never forgets.
\end{itemize}

FIR systems are computed directly as a finite convolution sum. IIR systems are
usually \emph{not} computed from $h$ (which is infinite!) but from a compact
recursive rule -- a difference equation, the subject of
Section~\ref{ch12:diffeq}.

\section{Causality and stability}
\label{ch12:causality-stability}

Two further properties decide whether a system can be \emph{built} and whether
it can be \emph{trusted}. They translate into clean conditions on $h[n]$.

\subsection*{Causality: no peeking at the future}

\begin{definition}[Causality]
\label{ch12:def-causal}
A system is \textbf{causal} if the output at time $n$ depends only on the input
at times $n$ and earlier -- never on future inputs. For an LTI system this is
equivalent to
\[
  h[n] = 0 \quad\text{for all } n < 0.
\]
\end{definition}

\begin{whybox}
Why does causality reduce to $h[n]=0$ for $n<0$? In the convolution sum
$y[n] = \sum_k x[k]\,h[n-k]$, a future input $x[k]$ with $k>n$ contributes only
through $h[n-k]$ where $n-k<0$. If $h$ is zero on negative indices, those future
terms vanish, so $y[n]$ uses only present and past inputs. Conversely, if
$h[n_*]\ne0$ for some $n_*<0$, then poking the input at a future time changes a
present output -- non-causal.
\end{whybox}

Any system that must run in \emph{real time} -- the firmware on the wrist
band reacting to motion as it happens -- \emph{must} be causal; it physically
cannot read samples that have not arrived. Non-causal systems are not useless,
though: when you process a recorded file offline, ``the future'' is already on
disk, and symmetric (non-causal) filters with zero phase distortion become
available. We exploit exactly this in Chapter~17.

\begin{trapbox}
A symmetric smoother like $y[n]=\tfrac13(x[n-1]+x[n]+x[n+1])$ is \emph{not}
causal -- it reads $x[n+1]$, one sample into the future. You can still implement
it in real time, but only by accepting a one-sample \emph{delay}: compute and
emit $y[n]$ once $x[n+1]$ is in hand, i.e.\ output the value labeled $n$ at clock
tick $n{+}1$. Linear-phase filters buy their pristine phase response with exactly
this kind of bulk latency. On a gesture controller that latency is part of your
end-to-end ``time from motion to MIDI'' budget.
\end{trapbox}

\subsection*{Stability: bounded in, bounded out (BIBO)}

\begin{definition}[BIBO stability]
\label{ch12:def-bibo}
A system is \textbf{BIBO stable} if every bounded input produces a bounded
output: whenever $\abs{x[n]}\le B_x<\infty$ for all $n$, there exists
$B_y<\infty$ with $\abs{y[n]}\le B_y$ for all $n$.
\end{definition}

For an LTI system there is a beautifully simple test in terms of $h$.

\begin{theorem}[Stability condition]
\label{ch12:thm-bibo}
An LTI system is BIBO stable if and only if its impulse response is
\emph{absolutely summable}:
\[
  \sum_{n=-\infty}^{\infty} \abs{h[n]} \;<\; \infty.
\]
\end{theorem}

\begin{whybox}
\textbf{($\Leftarrow$ sufficiency)} Suppose $S=\sum_n\abs{h[n]}<\infty$ and
$\abs{x[n]}\le B_x$. Bound the output directly:
\[
  \abs{y[n]} = \left|\sum_k h[k]\,x[n-k]\right|
  \le \sum_k \abs{h[k]}\,\abs{x[n-k]}
  \le B_x \sum_k \abs{h[k]} = B_x\,S < \infty.
\]
So the output is bounded by $B_x S$ for every $n$ -- stable.

\textbf{($\Rightarrow$ necessity)} If $S=\infty$, we construct a bounded input
that blows up. Pick the bounded input $x[n] = \operatorname{sign} h[-n]$ (which
satisfies $\abs{x[n]}\le1$). Then at $n=0$,
$y[0]=\sum_k h[k]\,x[-k]=\sum_k h[k]\operatorname{sign}h[k]=\sum_k\abs{h[k]}=
\infty$. A bounded input produced an unbounded output -- unstable.
\end{whybox}

\begin{example}[Which of our systems are stable?]
\begin{itemize}
  \item Averager $h=(\tfrac12,\tfrac12)$: $\sum\abs{h}=1<\infty$. Stable.
  \item Accumulator $h[n]=u[n]$: $\sum_{n\ge0}1=\infty$. \emph{Unstable} -- feed
        it a constant $x[n]=1$ and the running sum grows without bound. (True of
        any pure integrator; in practice we tame it with feedback or leakage.)
  \item Decaying echo $h[n]=a^n u[n]$ with $\abs a<1$: geometric series
        $\sum_{n\ge0}\abs a^n = \tfrac{1}{1-\abs a}<\infty$. Stable. With
        $\abs a\ge1$ it diverges -- unstable.
\end{itemize}
\end{example}

\begin{connectionbox}
The geometric-series threshold $\abs a<1$ is not a coincidence. In Chapter~15 we
will see that $a$ is a \emph{pole} of the system's $z$-transform, and the
stability condition ``$\abs a<1$'' becomes the elegant geometric statement
\emph{all poles lie strictly inside the unit circle}. Stability, which we just
met as a summability inequality, will turn into a picture you can read at a
glance from a pole--zero plot.
\end{connectionbox}

\section{Difference equations: how systems are really run}
\label{ch12:diffeq}

Convolution describes \emph{what} an LTI system does. But you rarely \emph{run}
an IIR system by convolving with its infinite $h$. Instead you run a
\textbf{linear constant-coefficient difference equation} (LCCDE): a finite recipe
relating a few output samples to a few input samples.

\begin{definition}[LCCDE]
\label{ch12:def-lccde}
A linear constant-coefficient difference equation has the form
\[
  \sum_{i=0}^{N} a_i\,y[n-i] \;=\; \sum_{j=0}^{M} b_j\,x[n-j],
\]
with constant coefficients $a_i, b_j$ and $a_0\ne0$. Solving for the current
output gives the directly runnable recursion
\[
  y[n] = \frac{1}{a_0}\left( \sum_{j=0}^{M} b_j\,x[n-j]
                            - \sum_{i=1}^{N} a_i\,y[n-i] \right).
\]
\end{definition}

The right-hand side is a \emph{recipe you can type into a loop}: to get the next
output, take a weighted blend of recent inputs (the $b_j$ part, the
\emph{feedforward} or FIR part) and subtract a weighted blend of recent outputs
(the $a_i$ part, the \emph{feedback} or recursive part). The feedback is what
makes the impulse response infinite from a finite rule.

\begin{example}[A one-pole smoother: leaky integrator]
\label{ch12:leaky}
Take $N=1$, $M=0$, $a_0=1$, $a_1=-a$, $b_0=(1-a)$:
\[
  y[n] = (1-a)\,x[n] + a\,y[n-1], \qquad 0<a<1.
\]
This is the classic exponential smoother (the digital first-order low-pass).
Its impulse response: feed $x=\delta$, assume rest ($y[n]=0$ for $n<0$):
\[
  y[0]=(1-a),\quad y[1]=a(1-a),\quad y[2]=a^2(1-a),\ \dots\ \Rightarrow\
  h[n]=(1-a)\,a^n u[n].
\]
A one-line recursion produces an \emph{infinite} geometric impulse response
(Figure~\ref{ch12:fig-leaky-impulse}).
That is the power -- and the danger -- of feedback: it is stable here only
because $\abs a<1$ makes the geometric series converge (cf.\
Theorem~\ref{ch12:thm-bibo}).
\end{example}

\begin{figure}[H]
\centering
\begin{tikzpicture}
\begin{axis}[
  width=11cm, height=4.2cm,
  title={Impulse response of the leaky integrator, $a=0.8$},
  title style={font=\small},
  xlabel={$n$}, ylabel={$h[n]$},
  xmin=-1.5, xmax=12.5, ymin=-0.03, ymax=0.26,
  xtick={0,2,4,6,8,10,12}, ytick={0,0.1,0.2},
  axis lines=left, ymajorgrids, grid style={gray!20},
]
\addplot+[ycomb, mark=*, sigblue, mark options={fill=sigblue}]
  coordinates {(-1,0)(0,0.2)(1,0.16)(2,0.128)(3,0.1024)(4,0.0819)(5,0.0655)
               (6,0.0524)(7,0.0419)(8,0.0336)(9,0.0268)(10,0.0215)(11,0.0172)(12,0.0137)};
\end{axis}
\end{tikzpicture}
\caption{The one-line recursion $y[n]=(1-a)x[n]+a\,y[n-1]$ produces the
geometric impulse response $h[n]=(1-a)a^n u[n]$ (here $a=0.8$, so
$h[0]=0.2$). It decays but never reaches zero -- an infinite impulse
response from finite memory.}
\label{ch12:fig-leaky-impulse}
\end{figure}

\begin{lstlisting}[language=Python,caption={The leaky integrator in a loop -- the cheapest useful real-time filter.}]
def leaky_integrator(x, a):
    y = 0.0
    out = []
    for xn in x:                 # one multiply-add per sample
        y = (1 - a) * xn + a * y # blend new input with smoothed history
        out.append(y)
    return out
# a near 1.0 -> heavy smoothing (long memory); a near 0 -> barely smooths.
\end{lstlisting}

\begin{focusbox}
\textbf{FIR vs.\ IIR, from the difference equation.} If there is no feedback
($N=0$: no $y[n-i]$ terms on the right), the system is FIR -- the output is a
finite window of inputs and $h$ is simply the list of $b_j$ coefficients. If
there \emph{is} feedback ($N\ge1$), the system is generally IIR. FIR is always
stable (a finite sum of bounded terms) and can have exactly linear phase; IIR is
cheaper (fewer taps for the same sharpness) but can go unstable and distorts
phase. We will weigh this trade-off carefully in Chapter~16.
\end{focusbox}

\section{Complex exponentials are the eigenfunctions of LTI systems}
\label{ch12:eigen}

We now reach the most beautiful fact in the chapter, the one that connects LTI
systems to everything Fourier. Feed an LTI system not an impulse but an
everlasting complex exponential, $x[n]=e^{\jj\omega n}$, and ask what comes out.

\begin{theorem}[Eigenfunction property]
\label{ch12:thm-eigen}
Let $T$ be LTI with impulse response $h[n]$. For the input $x[n]=e^{\jj\omega n}$,
\[
  y[n] \;=\; H(e^{\jj\omega})\,e^{\jj\omega n},
  \qquad\text{where}\qquad
  H(e^{\jj\omega}) \;=\; \sum_{k=-\infty}^{\infty} h[k]\,e^{-\jj\omega k}.
\]
The output is the \emph{same} complex exponential, merely multiplied by the
complex number $H(e^{\jj\omega})$.
\end{theorem}

\begin{figure}[H]
\centering
\begin{tikzpicture}[
  >=Stealth,
  sys/.style={draw=sourceblue!80, fill=sourceblue!10, rounded corners=2pt, very thick, minimum width=36mm, minimum height=14mm, align=center},
  sig/.style={font=\small},
]
\node[sig] (in) {$x[n] = e^{\jj\omega n}$};
\node[sys, right=14mm of in] (box) {\textbf{LTI System}\\[0.2em]Impulse response $h[n]$};
\node[sig, right=14mm of box] (out) {$y[n] = H(e^{\jj\omega})\,e^{\jj\omega n}$};
\draw[->, very thick, sigblue] (in) -- (box);
\draw[->, very thick, sigblue] (box) -- (out);

\node[below=3mm of in, font=\footnotesize, text=gray!80!black, align=center] {Pure complex phasor\\at frequency $\omega$};
\node[below=3mm of out, font=\footnotesize, text=specred, align=center] {Same frequency $\omega$, scaled by\\eigenvalue $H(e^{\jj\omega}) = |H|e^{\jj\angle H}$};
\node[above=3mm of box, font=\footnotesize\scshape, text=sourceblue!80!black] {Eigenfunction Property};
\end{tikzpicture}
\caption{Complex exponentials are the eigenfunctions of LTI systems. The system cannot create new frequencies; it can only scale the magnitude by $|H(e^{\jj\omega})|$ and shift the phase by $\angle H(e^{\jj\omega})$.}
\label{ch12:fig-eigen-block}
\end{figure}

\begin{whybox}
Plug $x[n]=e^{\jj\omega n}$ into the convolution sum and factor:
\[
  y[n] = \sum_{k} h[k]\,x[n-k]
       = \sum_{k} h[k]\,e^{\jj\omega(n-k)}
       = e^{\jj\omega n}\underbrace{\sum_{k} h[k]\,e^{-\jj\omega k}}_{H(e^{\jj\omega})}.
\]
The shift-invariance of the exponential is the whole trick: $e^{\jj\omega(n-k)} =
e^{\jj\omega n}\,e^{-\jj\omega k}$ splits into the part that depends on $n$
(pulled out front) and the part that depends on $k$ (summed against $h$). The
$k$-sum no longer depends on $n$, so it is just a constant scaling.
\end{whybox}

\begin{happeningbox}
This is exactly the eigenvector idea from linear algebra. A matrix $\mat A$
usually rotates and stretches a vector into a new direction; but along an
\emph{eigenvector} it merely scales, $\mat A\vec v=\lambda\vec v$. Here the
``matrix'' is the LTI system $T$, the ``eigenvectors'' are the complex
exponentials $e^{\jj\omega n}$, and the ``eigenvalues'' are the complex numbers
$H(e^{\jj\omega})$. \emph{Every LTI system is diagonal in the basis of complex
exponentials.} That is the deep reason the Fourier transform is the natural
language of filtering: it changes coordinates into the basis where every LTI
system is just a list of scalings.
\end{happeningbox}

\begin{connectionbox}
The function $H(e^{\jj\omega})=\sum_k h[k]e^{-\jj\omega k}$ is precisely the
\textbf{discrete-time Fourier transform} of the impulse response (Chapter~8). So
two pictures of the same system meet here: the \emph{time-domain} picture is the
impulse response $h[n]$, and the \emph{frequency-domain} picture is its DTFT
$H(e^{\jj\omega})$. Convolution in time (Theorem~\ref{ch12:thm-conv}) becomes
\emph{multiplication} in frequency, $Y(e^{\jj\omega})=H(e^{\jj\omega})\,
X(e^{\jj\omega})$ -- the convolution theorem you met for signals now reread as a
statement about systems.
\end{connectionbox}

\subsection*{The frequency response: magnitude and phase}

The eigenvalue $H(e^{\jj\omega})$ is the \textbf{frequency response}. Being
complex, it carries two real pieces of information at each frequency $\omega$:
\[
  H(e^{\jj\omega}) = \abs{H(e^{\jj\omega})}\,e^{\jj\angle H(e^{\jj\omega})}.
\]
\begin{itemize}
  \item $\abs{H(e^{\jj\omega})}$ -- the \textbf{magnitude response} -- says how
        much the system amplifies ($>1$) or attenuates ($<1$) a pure tone at
        frequency $\omega$.
  \item $\angle H(e^{\jj\omega})$ -- the \textbf{phase response} -- says how much
        the system delays (shifts) that tone.
\end{itemize}
Because real sinusoids are sums of conjugate exponentials,
$\cos(\omega n)=\tfrac12(e^{\jj\omega n}+e^{-\jj\omega n})$, the eigenfunction
property gives the practical rule: a real LTI system turns the input
$\cos(\omega n)$ into
\[
  y[n] = \abs{H(e^{\jj\omega})}\,\cos\!\big(\omega n + \angle H(e^{\jj\omega})\big).
\]
\emph{A sinusoid in gives the same-frequency sinusoid out, rescaled in amplitude
and shifted in phase -- never a new frequency.} This ``no new frequencies''
property is the signature of LTI systems and the reason nonlinear systems (the
squarer, the rectifier) are so different: they manufacture harmonics.

\begin{figure}[H]
\centering
\begin{tikzpicture}
\begin{axis}[
  width=12cm, height=5cm,
  title={Sinusoid in, same-frequency sinusoid out (scaled and phase-shifted)},
  title style={font=\small},
  xlabel={$n$}, xmin=0, xmax=20, ymin=-1.3, ymax=1.3,
  samples=120, domain=0:20,
  xtick={0,5,10,15,20}, ytick={-1,0,1},
  axis lines=left, ymajorgrids, grid style={gray!20},
  legend pos=south east, legend style={font=\footnotesize, draw=gray!40},
]
\addplot[sigblue, very thick] {cos(deg(0.5*x))};
\addlegendentry{input $\cos(\omega n)$}
\addplot[specred, very thick, dashed] {0.6*cos(deg(0.5*x - 0.9))};
\addlegendentry{output $|H|\cos(\omega n + \angle H)$}
\end{axis}
\end{tikzpicture}
\caption{An LTI system maps a pure tone to a tone at the \emph{same} frequency,
multiplied by the magnitude $\abs{H(e^{\jj\omega})}$ (here a gain of $0.6$) and
slid by the phase $\angle H(e^{\jj\omega})$ (here a lag of about one sample).
The shape and frequency are untouched.}
\label{ch12:fig-eigen-sinusoid}
\end{figure}

\begin{example}[Frequency response of the two-point averager]
\label{ch12:freq-averager}
For $h=(\tfrac12,\tfrac12)$,
\[
  H(e^{\jj\omega}) = \tfrac12 + \tfrac12 e^{-\jj\omega}
  = e^{-\jj\omega/2}\,\cos(\omega/2),
\]
where the last step factors out the average phase. Hence
$\abs{H(e^{\jj\omega})}=\abs{\cos(\omega/2)}$ and the phase is the pure linear
delay $-\omega/2$. At DC ($\omega=0$) the magnitude is $1$ -- constants pass
untouched. At the highest digital frequency ($\omega=\pi$, the fastest
alternation $\dots,+1,-1,+1,\dots$) the magnitude is $\cos(\pi/2)=0$ -- the
averager kills it completely. So the two-point averager is a (gentle) low-pass
filter, exactly as intuition demanded.
\end{example}

\begin{figure}[H]
\centering
\begin{tikzpicture}
\begin{axis}[
  width=11.5cm, height=5cm,
  title={Magnitude response of the two-point averager},
  title style={font=\small},
  xlabel={$\omega$ (rad/sample)}, ylabel={$|H(e^{\mathrm{j}\omega})|$},
  domain=0:3.14159, samples=120,
  xmin=0, xmax=3.14159, ymin=-0.05, ymax=1.1,
  xtick={0,0.7854,1.5708,2.3562,3.14159},
  xticklabels={$0$,$\pi/4$,$\pi/2$,$3\pi/4$,$\pi$},
  ytick={0,0.5,1}, axis lines=left, ymajorgrids, grid style={gray!20},
]
\addplot[specred, very thick] {abs(cos(deg(x/2)))};
\addplot[gray, dashed, thick] coordinates {(0,0.7071)(1.5708,0.7071)};
\addplot[gray, dashed, thick] coordinates {(1.5708,0)(1.5708,0.7071)};
\node[anchor=south west, font=\footnotesize, gray!50!black]
  at (axis cs:0.1,0.73) {half-power $\approx 0.707$ at $\omega = \pi/2$};
\end{axis}
\end{tikzpicture}
\caption{The averager passes low frequencies (gain $1$ at $\omega=0$) and
rejects the fastest alternation (gain $0$ at $\omega=\pi$): a simple low-pass.}
\label{ch12:fig-averager-mag}
\end{figure}

The mirror image is the \emph{first-difference} system $h=(1,-1)$, the simplest
edge detector. Its magnitude response $\abs{H(e^{\jj\omega})}=2\abs{\sin(\omega/2)}$
is zero at DC and largest at $\omega=\pi$: a high-pass that ignores slow trends
and reacts to fast change (Figure~\ref{ch12:fig-diff-mag}). You will derive it in
the exercises.

\begin{figure}[H]
\centering
\begin{tikzpicture}
\begin{axis}[
  width=11.5cm, height=5cm,
  title={Magnitude response of the first-difference (edge detector)},
  title style={font=\small},
  xlabel={$\omega$ (rad/sample)}, ylabel={$|H(e^{\mathrm{j}\omega})|$},
  domain=0:3.14159, samples=120,
  xmin=0, xmax=3.14159, ymin=-0.1, ymax=2.2,
  xtick={0,0.7854,1.5708,2.3562,3.14159},
  xticklabels={$0$,$\pi/4$,$\pi/2$,$3\pi/4$,$\pi$},
  ytick={0,1,2}, axis lines=left, ymajorgrids, grid style={gray!20},
]
\addplot[specred, very thick] {2*abs(sin(deg(x/2)))};
\end{axis}
\end{tikzpicture}
\caption{The first difference $h=(1,-1)$ rejects DC (gain $0$ at $\omega=0$) and
peaks at the fastest alternation (gain $2$ at $\omega=\pi$): a high-pass,
the complement of the averager.}
\label{ch12:fig-diff-mag}
\end{figure}

\begin{aestheticbox}
Notice how much rests on a single sequence. \textbf{One} object, the impulse response
$h[n]$, simultaneously:
(i) generates every output by convolution;
(ii) decides causality by where it is nonzero;
(iii) decides stability by whether it is absolutely summable;
(iv) \emph{is} the system's frequency response once you take its Fourier
transform. Time-domain shape and frequency-domain behavior are two readings of a
single page. Mastering LTI systems means becoming fluent in flipping between
those two readings at will.
\end{aestheticbox}

\section{Block diagrams: drawing systems}
\label{ch12:block}

Engineers rarely write difference equations on the workbench; they \emph{draw}
them. A block diagram is a wiring picture built from three primitives:
\begin{itemize}
  \item an \textbf{adder} (a circle with $+$) that sums incoming signals;
  \item a \textbf{gain} (a triangle or labeled arrow) that scales by a constant;
  \item a \textbf{unit delay} (a box labeled $z^{-1}$) that holds a sample for
        one clock tick -- this is the hardware register that gives the system its
        memory.
\end{itemize}
The label $z^{-1}$ is borrowed from the $z$-transform (Chapter~15), where
delaying by one sample corresponds to multiplying by $z^{-1}$. For now read it
simply as ``delay by one.''

\begin{figure}[H]
\centering
\begin{tikzpicture}[
  >=Stealth,
  blk/.style={draw=sourceblue, rounded corners=2pt, minimum width=10mm, minimum height=7mm, fill=sourceblue!10, font=\small},
  sum/.style={draw=targetgreen!80!black, circle, inner sep=0pt, minimum size=7mm, fill=targetgreen!15, font=\normalsize},
  gain/.style={draw=quotientorange!80!black, regular polygon, regular polygon sides=3, shape border rotate=-90,
               inner sep=0pt, minimum size=12mm, fill=quotientorange!15},
  gainleft/.style={draw=quotientorange!80!black, regular polygon, regular polygon sides=3, shape border rotate=90,
               inner sep=0pt, minimum size=11mm, fill=quotientorange!15},
  dot/.style={circle, fill=black, inner sep=1.2pt},
]
% Forward path
\node (x) at (0,0) {$x[n]$};
\node[gain] (g1) at (1.15,0) {};
\node[font=\footnotesize] at (g1.center) {$1{-}a$};

\node[sum] (s) at (2.5,0) {$+$};
\node[dot] (tap) at (3.9,0) {};
\node (y) at (5.1,0) {$y[n]$};

% Forward wires
\draw[->, thick] (x) -- (g1.west);
\draw[->, thick] (g1.east) -- (s.west);
\draw[thick] (s) -- (tap);
\draw[->, thick] (tap) -- (y);

% Feedback path: delay under tap, gain under forward path
\node[blk] (z) at (3.9,-1.5) {$z^{-1}$};
\node[gainleft] (g2) at (3.0,-1.5) {};
\node[font=\footnotesize] at (g2.center) {$a$};

% Feedback wires
\draw[->, thick] (tap) -- (z.north);
\draw[->, thick] (z.west) -- (g2.east);
\draw[->, thick] (g2.west) -| (s.south);

% Annotations
\node[font=\tiny, text=gray!80!black, anchor=west] at (4.5,-1.5) {delay};
\node[font=\tiny, text=quotientorange!90!black, anchor=north] at (3.0,-2.0) {feedback};
\end{tikzpicture}
\caption{Block diagram of the leaky integrator
$y[n]=(1-a)x[n]+a\,y[n-1]$. The $z^{-1}$ block stores the previous output; the
loop back through gain $a$ is the feedback that makes the impulse response
infinite.}
\label{ch12:fig-block-leaky}
\end{figure}

\begin{walkbox}
\textbf{Read a block diagram back into an equation.} Trace from the output adder
backward. In Figure~\ref{ch12:fig-block-leaky} the adder produces $y[n]$ by
summing two paths: the input scaled by $(1-a)$, and the previous output (the
$z^{-1}$ box turns $y[n]$ into $y[n-1]$) scaled by $a$. Writing that sum gives
$y[n]=(1-a)x[n]+a\,y[n-1]$ -- back to Example~\ref{ch12:leaky}. Diagram,
equation, impulse response, and frequency response are four faces of one system;
fluency means moving among them without friction.
\end{walkbox}

\section{Why this matters for the wrist controller}
\label{ch12:wrist}

\begin{goalbox}
Almost every processing block in the gesture controller is, by design, LTI -- and
we keep it that way on purpose, because LTI is the part we can reason about with
the frequency response.

\textbf{The LTI parts.} The IMU gives accelerometer and gyroscope streams riddled
with high-frequency vibration and electrical noise. The first thing we do is
\emph{low-pass} each axis -- often with the very leaky integrator of
Example~\ref{ch12:leaky}, one multiply-add per sample, needing minimal computation on an
8-bit micro. Removing the gravity component to isolate dynamic motion is a
\emph{high-pass}; isolating a forearm-tremor band is a \emph{band-pass}. All
three are LTI filters, fully described by their impulse responses, fully
predictable from their frequency responses.

\textbf{The deliberately non-LTI parts.} Gesture recognition itself is \emph{not}
LTI, and that is the point. To turn EMG into a ``pinch detected'' event we
\emph{rectify} ($y=\abs{x}$, nonlinear -- Example~\ref{ch12:squarer-nonlinear}'s
cousin) and then threshold (also nonlinear). Nonlinearity is exactly what lets us
collapse a continuous waveform into a discrete decision. The clean architecture
is: \emph{LTI filtering first} (where the math is friendly), \emph{nonlinear
decision last} (where we actually want new behavior, like emitting a MIDI
note-on).

\textbf{Causality is non-negotiable here.} The controller must react as motion
happens; it cannot read future samples. Every filter in the live path must be
causal ($h[n]=0$ for $n<0$), and every sample of group delay it adds is latency
between your gesture and the sound. Budget it.
\end{goalbox}

\begin{trapbox}
\textbf{What the sensors can and cannot tell you.} It is tempting to imagine an LTI filter bank
that ``decodes'' individual finger joint angles from a single wrist IMU. It
cannot, and no amount of filtering changes that. A wrist-mounted IMU measures the
motion of the \emph{wrist}; it has no observability of independent finger flexion
when the wrist is still. The most an IMU-only band can do reliably is detect
\emph{coarse, discrete} gestures -- a flick, a tap, a wrist rotation, a shake.
Recovering finger-level activity requires sensing closer to the muscles or
fingers: surface EMG or force-myography (FMG) on the forearm for discrete
finger gestures, or instrumented rings / data gloves for continuous finger pose.
This is consistent with the published gesture-sensing literature (e.g.\ Myo-style
forearm EMG armbands for discrete finger gestures, and FMG cuffs for grip
states). LTI filtering cleans up whatever the sensor can observe; it cannot
manufacture information the sensor never captured.
\end{trapbox}

\begin{applicationbox}
The same LTI ideas reappear everywhere in this book and in the wild:
\emph{audio} -- a graphic equalizer is a parallel bank of LTI band-pass filters,
and a convolution reverb literally convolves your signal with the measured
impulse response of a concert hall; \emph{radio} -- the channel between
transmitter and receiver is modeled as an LTI system whose impulse response is
the multipath profile; \emph{biomedical} -- the 50/60\,Hz mains hum in an ECG is
removed by an LTI notch filter; \emph{seismology} -- the earth's response to a
distant quake is the convolution of the source with the LTI transfer function of
the rock between. Learn LTI once, recognize it forever.
\end{applicationbox}

\section{A continuous-time footnote}
\label{ch12:ct}

Everything above has a continuous-time mirror. A continuous LTI system is
described by its impulse response $h(t)=T\{\delta(t)\}$, the output is the
convolution \emph{integral}
\[
  y(t) = (x\conv h)(t) = \int_{-\infty}^{\infty} x(\tau)\,h(t-\tau)\,\dd\tau,
\]
the eigenfunctions are $e^{\jj\Omega t}$ with eigenvalue
$H(\jj\Omega)=\int h(t)e^{-\jj\Omega t}\dd t$ (the continuous Fourier transform of
$h$), causality is $h(t)=0$ for $t<0$, and BIBO stability is
$\int\abs{h(t)}\dd t<\infty$. The analog smoothing $RC$ circuit in front of your
ADC is a continuous-time first-order low-pass; its discrete cousin is the leaky
integrator. Same geometry, integrals instead of sums.

\section{Summary}
\label{ch12:summary}

\begin{recoverybox}
\textbf{Key takeaways.}
\begin{itemize}
  \item A \emph{system} maps input signals to output signals.
  \item \emph{Linear} = respects scaling and addition (superposition);
        \emph{time-invariant} = no clock, shifting the input just shifts the
        output. \emph{LTI} = both.
  \item Any signal is a sum of scaled, shifted impulses (sifting). Apply
        linearity then time invariance and you get: \emph{LTI output = input
        convolved with the impulse response} $h[n]$
        (Theorem~\ref{ch12:thm-conv}). One measured $h$ predicts all behavior.
  \item \emph{Causal} $\Leftrightarrow h[n]=0$ for $n<0$ (required for real time).
        \emph{BIBO stable} $\Leftrightarrow \sum_n\abs{h[n]}<\infty$.
  \item FIR (finite $h$, no feedback, always stable, linear phase possible) vs.\
        IIR (infinite $h$, feedback, cheaper, can be unstable). Both run as
        difference equations.
  \item Complex exponentials are the \emph{eigenfunctions} of LTI systems; the
        eigenvalue is the \emph{frequency response} $H(e^{\jj\omega})$, the DTFT
        of $h$. A sinusoid in gives the same-frequency sinusoid out, scaled by
        $\abs{H}$ and phase-shifted by $\angle H$ -- never a new frequency.
  \item Block diagrams ($+$, gain, $z^{-1}$) are the engineer's drawing of the
        same difference equation.
\end{itemize}
\end{recoverybox}

\begin{connectionbox}
\textbf{Capabilities, Downsides, and Forward Roadmap}
\begin{itemize}
  \item \textbf{What this unlocked}: The LTI framework establishes that an entire system is governed by
  a single signal: its impulse response $h[n]$. Once we measure or specify $h[n]$, superposition and time invariance
  uniquely dictate the output for any arbitrary input.
  \item \textbf{Honest downsides and limitations}: While the convolution sum $y[n] = \sum_k x[k]h[n-k]$
  was introduced here algebraically, its mechanical execution hides practical hurdles: how to compute it
  efficiently, how boundary conditions and edge transients unfold, and how to avoid the quadratic runtime
  $\mathcal{O}(NM)$ for long signals.
  \item \textbf{Where we go next}: Chapter~13 turns to \textbf{Convolution}: we work through the
  flip-and-slide mechanics, the link to polynomial multiplication, fast convolution via the FFT,
  and 2D convolution for images.
\end{itemize}
\end{connectionbox}

\section{Exercises}
\label{ch12:exercises}

\begin{exercise}[Classify the zoo]
For each system in Example~\ref{ch12:zoo}, state whether it is (a) linear and
(b) time-invariant, with a one-line justification or counterexample for each.
Which are LTI?
\end{exercise}

\begin{exercise}[Affine is not linear]
Show that $y[n]=2x[n]-1$ violates Definition~\ref{ch12:def-linear} by exhibiting
two signals $x_1,x_2$ for which $T\{x_1+x_2\}\ne T\{x_1\}+T\{x_2\}$. Then explain
in one sentence why \emph{any} nonzero constant term ruins linearity.
\end{exercise}

\begin{exercise}[Time-variant detector]
Consider $y[n]=x[n]\cos(\omega_0 n)$ (amplitude modulation by a fixed carrier).
Prove it is linear but \emph{not} time-invariant. Where does the free $n$ enter?
\end{exercise}

\begin{exercise}[Convolve by hand]
Let $h=(1,-1)$ at $n=0,1$ (the first-difference / edge detector) and
$x=(2,2,5,5,5)$ at $n=0,\dots,4$. Compute $y=x\conv h$ by marching $n$ forward
as in Example~\ref{ch12:hand-conv}. Where is the output largest, and why does
that make sense for an edge detector?
\end{exercise}

\begin{exercise}[Impulse response of a difference equation]
A system obeys $y[n]=x[n]-x[n-3]$. Find $h[n]$ by feeding $x=\delta$. Is it FIR
or IIR? Causal? Stable (check $\sum\abs{h[n]}$)?
\end{exercise}

\begin{exercise}[Leaky-integrator impulse response]
For $y[n]=(1-a)x[n]+a\,y[n-1]$ with $a=0.9$, compute $h[0],h[1],h[2],h[3]$
numerically and confirm they follow $h[n]=(1-a)a^n$. What is $\sum_{n\ge0}h[n]$,
and why does that value mean ``unity DC gain''?
\end{exercise}

\begin{exercise}[Stability boundary]
For $h[n]=a^n u[n]$, the system is BIBO stable iff $\abs a<1$. (a) Evaluate
$\sum_n\abs{h[n]}$ for $a=0.5$ and for $a=1.5$. (b) For $a=1$ exactly (the
accumulator restricted to a single pole on the boundary), is the system stable?
Give a bounded input whose output is unbounded.
\end{exercise}

\begin{exercise}[Frequency response from $h$]
For the first-difference system $h=(1,-1)$, show that
$H(e^{\jj\omega})=1-e^{-\jj\omega}=2\jj\,e^{-\jj\omega/2}\sin(\omega/2)$. Sketch
$\abs{H(e^{\jj\omega})}=2\abs{\sin(\omega/2)}$ over $0\le\omega\le\pi$. Is it
low-pass or high-pass? Connect this to its role as an edge detector.
\end{exercise}

\begin{exercise}[Sinusoid in, sinusoid out]
The two-point averager (Example~\ref{ch12:freq-averager}) has
$H(e^{\jj\omega})=e^{-\jj\omega/2}\cos(\omega/2)$. If the input is
$x[n]=\cos(\tfrac{\pi}{2}n)$, write the exact output $y[n]$ including its
amplitude and phase shift. By how many samples is the tone delayed?
\end{exercise}

\begin{exercise}[Cascade of LTI systems]
Two LTI systems with impulse responses $h_1$ and $h_2$ are connected in series
(output of the first feeds the second). Using Theorem~\ref{ch12:thm-conv} twice,
show the overall impulse response is $h_1\conv h_2$, and argue from
$Y=H_2(H_1 X)$ that the overall frequency response is the \emph{product}
$H_1(e^{\jj\omega})H_2(e^{\jj\omega})$. Why does this prove that the \emph{order}
of two LTI blocks does not matter?
\end{exercise}

\begin{exercise}[Nonlinearity makes harmonics]
Feed $x[n]=\cos(\omega_0 n)$ into the squarer $y[n]=x[n]^2$. Using
$\cos^2\theta=\tfrac12+\tfrac12\cos(2\theta)$, show the output contains a DC term
and a tone at $2\omega_0$ -- a frequency \emph{not present} in the input.
Contrast this with the ``no new frequencies'' guarantee for LTI systems, and
explain why this property is precisely what makes a rectifier useful for turning
an oscillating EMG burst into a detectable envelope.
\end{exercise}

\begin{exercise}[Draw and read a block diagram]
Draw the block diagram (using adders, gains, and $z^{-1}$ blocks) for the FIR
filter $y[n]=0.25x[n]+0.5x[n-1]+0.25x[n-2]$. How many delay elements does it
need? Then read your diagram back into the difference equation to check it. Is
this filter's phase linear? (Hint: look at the symmetry of its coefficients.)
\end{exercise}
